% ============================================================
% methodology.tex
% ============================================================


This chapter formalizes the statistical framework used to model
time-to-resolution and competing outcomes in federal securities
class actions. The central research question is: \emph{how do
legal regime changes (the PSLRA), court geography (federal
circuit), and case characteristics jointly shape both the
timing and the type of resolution---settlement versus
dismissal---in federal securities class actions?} Answering
this requires estimating how quickly cases resolve, whether
they resolve via settlement or dismissal, and how observable
characteristics alter those probabilities. Because settlement
and dismissal are mutually exclusive terminal events, standard
single-event survival methods are insufficient, requiring a
competing-risks framework throughout.

\section{Notation and Data Structure}

Let $i = 1, \ldots, n$ index individual securities class action
cases. For each case, we observe:
\begin{itemize}
    \item $T_i \geq 0$: time from filing to resolution,
    measured in years.
    \item $\delta_i \in \{0, 1, 2\}$: event type indicator,
    where $\delta_i = 0$ denotes right-censoring (case still
    pending), $\delta_i = 1$ denotes settlement, and 
    $\delta_i = 2$ denotes dismissal.
    \item $\mathbf{X}_i = (X_{i1}, \ldots, X_{ip})^\top$: a 
    $p$-dimensional vector of covariates measured at or near
    case filing.
\end{itemize}

The observed data are $\{(T_i, \delta_i,
\mathbf{X}_i)\}_{i=1}^n$. The time scale is years since
filing. There is no left truncation. Right censoring at the
IDB snapshot date is assumed to be conditionally
non-informative: after conditioning on observed covariates,
especially filing year, the probability of censoring is
taken to be independent of the latent event time. In a
1990--2024 cohort this assumption is not innocuous, because
recently filed cases are more likely to be censored by
construction; conditioning on calendar time partly mitigates
but does not eliminate that concern.

\subsection{Covariate Definitions}

\textbf{Core covariates} (available for all cases):
\begin{itemize}
    \item $\text{PSLRA}_i \in \{0,1\}$: binary indicator equal
    to 1 if case $i$ was filed on or after December~22, 1995.
    \item $\text{Circuit}_i$: federal circuit of filing,
    categorical with the Second Circuit as reference.
    \item $\text{FilingYear}_i$: calendar year of filing,
    used in the spline robustness specification in
    Section~\ref{sec:temporal_id} rather than in the primary
    extended model.
\end{itemize}

\textbf{Extended covariates} (available for cases with
non-missing IDB fields):
\begin{itemize}
    \item $\text{Origin}_i \in \{\text{Original},
    \text{MDL}, \text{Other}\}$: how the case entered federal
    court, with ``Removed'' (cases removed from state court)
    collapsed into ``Other'' to avoid thin cells.
    ``Original'' is the reference category.
    \item $\text{MDLFlag}_i \in \{0,1\}$: binary indicator
    equal to 1 if the IDB \texttt{mdldock} field is populated,
    indicating MDL docket assignment.
    \item $\text{JurisFQ}_i \in \{0,1\}$: binary indicator
    equal to 1 if the case invokes federal question
    jurisdiction (IDB \texttt{juris} code~3).
    \item $\text{StatBasis}_i \in \{\text{10(b)},
    \text{Section~11}, \text{Other/Both}, \text{Missing}\}$:
    statutory basis coded from the IDB \texttt{section} field.
    Cases with missing statutory basis are coded as an explicit
    ``Missing'' factor level rather than dropped, preserving
    the full estimation sample.
\end{itemize}

\section{Nonparametric Estimation of Cumulative Incidence}

The \emph{overall survival function} $S(t) = P(T > t)$ estimates the probability that a case remains unresolved beyond time $t$. We estimate $S(t)$ using the Kaplan-Meier estimator:
\begin{equation}
\hat{S}(t) = \prod_{t_j \leq t} \left(1 - \frac{d_j}{n_j}\right)
\end{equation}
where $t_j$ are the distinct observed event times, $d_j$ is the total number of events at $t_j$, and $n_j$ is the number of cases at risk just before $t_j$ \citep{kaplan1958}.

The cumulative incidence function (CIF) for event type $k \in \{1,2\}$ is:
\begin{equation}
F_k(t) = P(T \leq t, \delta = k)
\end{equation}
which represents the probability that a case has experienced event type $k$ by time $t$. The Aalen-Johansen estimator is given by:
\begin{equation}
\hat{F}_k(t) = \sum_{t_j \leq t} \hat{S}(t_{j-1}) \frac{d_{kj}}{n_j}
\end{equation}
where $d_{kj}$ is the number of type-$k$ events at $t_j$ \citep{aalen1978}. We plot $\hat{F}_1(t)$ and $\hat{F}_2(t)$ stratified by PSLRA regime and circuit. Equality of CIFs across groups is tested using Gray's $K$-sample test \citep{gray1988}.

\section{Semi-Parametric Hazard Regression}
\label{sec:hazard_regression}

Three progressively more informative modeling frameworks
address different aspects of the research question.
Cause-specific Cox models estimate how covariates alter
the instantaneous rate of each outcome among cases still at
risk, but they cannot directly quantify how those rate
changes translate into cumulative outcome probabilities
because they hold the competing pathway fixed. Fine-Gray
subdistribution models fill this gap by directly
parameterizing the cumulative incidence function, revealing
how covariates shift the marginal probability of each
outcome after accounting for the competing event. Neither
framework, however, separates the PSLRA's legislative
association from concurrent shifts in case
composition---a distinction that requires the propensity-score
weighting developed in Section~\ref{sec:iptw}.

\subsection{Cause-Specific Cox Proportional Hazards Models}

The Cox model relates each case characteristic to its association
with the speed of resolution, holding all other factors constant.
The key output is the \emph{hazard ratio} (HR): for a binary
covariate such as PSLRA, a hazard ratio of 1.50 would mean that,
at any point in time, post-PSLRA cases face a 50\% higher
instantaneous rate of that event compared to pre-PSLRA cases,
conditional on other covariates. $\text{HR} > 1$ indicates an
accelerated event rate; $\text{HR} < 1$ indicates a suppressed
rate; and $\text{HR} = 1$ indicates no association.

Formally, the \emph{cause-specific hazard} for event type $k$ is defined as:
\begin{equation}
\lambda_k(t \mid \mathbf{X}) = \lim_{\Delta t \to 0} \frac{P(t \leq T < t + \Delta t, \delta = k \mid T \geq t, \mathbf{X})}{\Delta t}
\end{equation}
Under the Cox proportional hazards model:
\begin{equation}
\lambda_k(t \mid \mathbf{X}) = \lambda_{0k}(t) \exp(\boldsymbol{\beta}_k^\top \mathbf{X})
\end{equation}
where $\lambda_{0k}(t)$ is the unspecified baseline hazard for event type $k$ and $\boldsymbol{\beta}_k$ is the vector of log-hazard ratios \citep{cox1972}. We estimate separate cause-specific Cox models for settlement ($k = 1$) and dismissal ($k = 2$). In each model, cases experiencing the competing event are treated as \emph{statistically censored}. This means they are removed from the risk set at the time of the competing event and treated as unobserved for the event of interest thereafter, not that their outcome is literally unknown.

\subsection{Fine-Gray Subdistribution Hazard Models}

While the cause-specific Cox model asks how covariates alter the
instantaneous rate of one event among cases still at risk, the
Fine-Gray model asks how covariates shift the \emph{cumulative
probability} of that event over time, accounting for the fact
that cases experiencing the competing event are permanently
removed from candidacy. Put differently, the cause-specific
model asks how quickly pending cases reach an outcome,
whereas the Fine-Gray model asks what share of all cases
reaches that outcome over follow-up. This distinction matters when the
research question concerns the overall incidence of an outcome
rather than its momentary rate.

The Fine-Gray model parameterizes the \emph{subdistribution hazard} \citep{fine1999}:
\begin{equation}
\lambda_k^*(t \mid \mathbf{X}) = \lambda_{0k}^*(t) \exp(\boldsymbol{\gamma}_k^\top \mathbf{X})
\end{equation}
where the subdistribution hazard $\lambda_k^*(t)$ is defined via:
\begin{equation}
F_k(t) = 1 - \exp\left(-\int_0^t \lambda_k^*(u) \, du\right)
\end{equation}
and $\boldsymbol{\gamma}_k$ is the subdistribution log-hazard ratio vector. Unlike cause-specific hazards, the subdistribution hazard has a direct marginal interpretation: $\exp(\gamma_{kj}) > 1$ is associated with a higher cumulative incidence of event $k$ over follow-up, while $\exp(\gamma_{kj}) < 1$ is associated with a lower cumulative incidence. We fit Fine-Gray models using the \texttt{finegray} function in R's \texttt{survival} package.

\subsection{Extended Covariate Specifications}

The baseline models include only the PSLRA indicator. We
then estimate a circuit-augmented specification with PSLRA
and circuit fixed effects, followed by extended
specifications that add origin category, MDL flag,
jurisdictional basis, and statutory basis. The extended
model linear predictor is:
\begin{align}
\eta_i ={}& \beta_{\text{PSLRA}} \text{PSLRA}_i + \boldsymbol{\beta}_{\text{circ}}^\top \mathbf{C}_i + \boldsymbol{\beta}_{\text{orig}}^\top \mathbf{O}_i \notag \\
&+ \beta_{\text{MDL}} \text{MDLFlag}_i + \beta_{\text{FQ}} \text{JurisFQ}_i + \boldsymbol{\beta}_{\text{stat}}^\top \mathbf{S}_i
\end{align}
where $\mathbf{C}_i$, $\mathbf{O}_i$, and $\mathbf{S}_i$ are indicator vectors for circuit, origin category, and statutory basis (including a ``Missing'' level) respectively, and $\eta_i$ enters the hazard as $\lambda_k(t \mid \mathbf{X}_i) = \lambda_{0k}(t)\exp(\eta_i)$.

\subsection{PSLRA $\times$ Circuit Interaction}

To test whether the PSLRA's association is uniform across circuits, we augment the extended specification with PSLRA $\times$ circuit interaction terms:
\begin{align}
\eta_i ={}& \beta_{\text{PSLRA}} \text{PSLRA}_i + \boldsymbol{\beta}_{\text{circ}}^\top \mathbf{C}_i + \boldsymbol{\alpha}^\top (\text{PSLRA}_i \cdot \mathbf{C}_i) \notag \\
&+ \boldsymbol{\beta}_{\text{orig}}^\top \mathbf{O}_i + \beta_{\text{MDL}} \text{MDLFlag}_i + \beta_{\text{FQ}} \text{JurisFQ}_i
\end{align}
where $\boldsymbol{\alpha}$ captures circuit-specific deviations from the baseline PSLRA coefficient. Statutory basis is excluded from the interaction specification because the PSLRA $\times$ circuit $\times$ statutory basis cross-tabulation produces thin cells that preclude stable estimation. Accordingly, both the additive and interaction versions of this circuit-heterogeneity exercise are estimated on the common narrower covariate set without statutory basis. The interaction model is compared to the additive specification using a likelihood ratio test (LRT), so the test is informative about interaction terms within that shared specification rather than about the richer extended model with statutory basis. We restrict the analysis to the top-six circuits by case volume to ensure stable estimates.

\subsection{Proportional Hazards Diagnostics and Remediation}
\label{sec:ph_diag}

The Cox model requires that hazard ratios remain constant over time. We test this assumption using the Grambsch-Therneau test, which regresses scaled Schoenfeld residuals on a function of time; a significant result ($p < 0.05$) indicates that the hazard ratio for that covariate varies over the study period \citep{grambsch1994}.

Our primary remediation targets the focal PSLRA covariate.
For the PSLRA indicator, we fit a piecewise constant hazard
model by splitting each case's time axis at $t = 1$ and
$t = 2$ years:
\begin{equation}
\lambda_k(t \mid \text{PSLRA}, \mathbf{X}) = \lambda_{0k}(t) \exp \left( \sum_{m=1}^{3} \beta_{km} \left( \text{PSLRA} \cdot \mathbf{1}_{I_m}(t) \right) + \boldsymbol{\beta}_{\text{other}}^\top \mathbf{X} \right)
\end{equation}
where $I_1 = [0,1)$, $I_2 = [1,2)$, and $I_3 = [2,\infty)$.
The cutpoints at $t = 1$ and $t = 2$ years were chosen
heuristically to align with the typical procedural timeline
of a motion to dismiss in securities class actions: the first
year captures the pleading-stage gauntlet; the second year
spans the interlocutory period before full discovery. These
cutpoints are not data-derived and do not represent formal
structural break estimates. The piecewise specification is
applied symmetrically to both the settlement and dismissal
models to allow direct comparison of temporal patterns across
outcomes.

For the remaining covariates, we retain the standard
constant-HR specifications and interpret their coefficients
as time-averaged summaries when PH is rejected. Formal PH
diagnostics for both the cause-specific Cox and Fine-Gray
models are reported in Chapter~\ref{ch:results}. Throughout
the thesis, ``proportional hazards assumption violations''
refers to statistical failures of the constant-HR
assumption, not legal infractions.

\section{Robustness and Sensitivity Analyses}

To assess the sensitivity of our central findings to
disposition code misclassification, cohort composition, and
circuit-specific confounding, we implement three sets of
robustness checks.

\textbf{Alternative disposition coding.} We implement three
coding schemes: Scheme~A (primary), Scheme~B (Code~12
reclassified as settlement), and Scheme~C (Codes~12 and~5
reclassified as settlement), as defined in Section~\ref{sec:coding_schemes}. We
re-estimate baseline PSLRA Cox models under all three schemes
and compare hazard ratios.

\textbf{Temporal restriction.} We exclude cases filed after
December~31, 2020, to assess whether recently filed cases with
high censoring rates distort estimates.

\textbf{Circuit-specific sub-samples.} We re-estimate PSLRA
Cox models separately within the Second and Ninth Circuits to
assess whether pooled estimates mask circuit-level
heterogeneity.

\subsection{Temporal Identification and Placebo Validation}
\label{sec:temporal_id}

Because the PSLRA indicator is perfectly collinear with
calendar time, the baseline PSLRA hazard ratio may conflate
the reform's legislative effect with secular trends in
judicial behavior. A na\"ive linear control for
$\text{FilingYear}_i$ forces 34~years of non-linear judicial
drift into a single slope, absorbing the PSLRA step-change
and producing spurious results. We instead implement three
complementary identification strategies.

\textbf{Natural spline time-trend control.} We add a natural
cubic spline basis for filing year with three degrees of
freedom, $\text{ns}(\text{FilingYear}_i, df{=}3)$, to the
cause-specific Cox model alongside $\text{PSLRA}_i$. The
spline flexibly captures non-linear secular trends---such as
evolving judicial attitudes or shifts in filing
composition---while leaving a separate PSLRA indicator in the
model to capture any discrete step-change at the cutoff.
Because the spline basis and the PSLRA indicator remain
strongly associated by construction, we report the variance
inflation factor (VIF) for PSLRA in this specification to
quantify the degree of remaining collinearity.

\textbf{Clean-window analysis.} We restrict the sample to
cases filed between 1993 and 1998---a narrow window
surrounding the PSLRA's December~22, 1995 effective
date---and re-estimate the baseline PSLRA Cox model without
any time-trend controls. The narrow window minimizes secular
confounding while preserving sufficient sample size for
estimation.

\textbf{Placebo cutoff test.} As a falsification check, we
construct a placebo ``law'' at January~1, 1992---a date with
no legislative significance---and re-estimate the Cox model
on the pre-PSLRA sample only (1990--1995). A significant
placebo effect would indicate that the model detects
pre-existing trends, casting doubt on the PSLRA
identification; an insignificant placebo effect strengthens
the case that the estimated PSLRA hazard ratio reflects
the reform itself rather than secular drift.

\section{Composition Adjustment via Propensity-Score Weighting}
\label{sec:iptw}

The Cox models in Section~\ref{sec:hazard_regression} estimate associations between
covariates and hazard rates. However, the pre- and post-PSLRA
case populations may differ systematically in composition---for
example, in the distribution of filing circuits, statutory
bases, or case origins---and these compositional shifts could
confound the estimated PSLRA effect. To isolate the PSLRA
component from concurrent changes in the observable case mix,
we employ inverse probability of treatment weighting
(IPTW) \citep{xieliu2005, austinfine2025}.

\subsection{Propensity Score Estimation}

The intuition is straightforward: we reweight pre-PSLRA cases
so that their observable characteristics---circuit, origin,
statutory basis---resemble those of post-PSLRA cases, isolating
the PSLRA component from concurrent compositional shifts.

Let $Z_i \in \{0,1\}$ denote the PSLRA indicator. The
propensity score is:
\begin{equation}
e(\mathbf{X}_i) = P(Z_i = 1 \mid \mathbf{X}_i)
\end{equation}
estimated via logistic regression of $Z_i$ on observed case
characteristics (circuit, origin, jurisdictional basis, and
statutory basis). Because the PSLRA is a legislative date
cutoff---not a treatment influenced by case characteristics---the
propensity score models \emph{temporal changes in case
composition} across eras, not a selection mechanism into
treatment.

\subsection{ATT Weighting}

We target the average treatment effect on the treated (ATT)
estimand, which re-weights the smaller pre-PSLRA control group
to match the covariate distribution of post-PSLRA cases. The
ATT weights are:
\begin{equation}
w_i =
\begin{cases}
1 & \text{if } Z_i = 1 \text{ (post-PSLRA)} \\[4pt]
\dfrac{\hat{e}(\mathbf{X}_i)}{1 - \hat{e}(\mathbf{X}_i)} & \text{if } Z_i = 0 \text{ (pre-PSLRA)}
\end{cases}
\end{equation}
ATT is preferred over the average treatment effect (ATE)
because the 11.6:1 post-to-pre sample imbalance would produce
extreme ATE weights. We trim weights at the 99th percentile
to limit the influence of extreme propensity scores and report
the effective sample size (ESS) as a measure of information
loss from weighting.

\subsection{Weighted Cox Models and Diagnostics}

The composition-adjusted hazard ratio is obtained from a
weighted cause-specific Cox model:
\begin{equation}
\lambda_k(t \mid Z_i) = \lambda_{0k}(t) \exp(\beta_k^{\text{w}} Z_i)
\end{equation}
where each observation contributes to the partial likelihood
with weight $w_i$ and robust (sandwich) standard errors account
for the estimated weights \citep{xieliu2005}. We implement a
four-strategy triangulation:
\begin{enumerate}
    \item \textbf{Unadjusted}: PSLRA only, no covariates, no weights.
    \item \textbf{Regression-adjusted}: Extended Cox (PSLRA + all
    covariates), no weights.
    \item \textbf{Covariate-Adjusted IPTW (Weighted Regression)}: Extended Cox with IPTW weights.
    \item \textbf{Marginal Structural Model}: PSLRA only, with IPTW weights.
\end{enumerate}
Agreement across these strategies is treated as a limited
sensitivity check on functional form rather than as evidence
from four fully independent estimators: strategies 2 and~3
share the same extended covariate set, while strategies 1
and~4 are distinguished primarily by weighting. Balance is assessed using standardized mean differences
(SMD), with $|\text{SMD}| < 0.1$
as the threshold for adequate balance \citep{austinfine2025}.

\subsection{Interpretive Scope}

We interpret IPTW-weighted hazard ratios as
\emph{composition-adjusted} estimates that isolate the PSLRA
component from observable shifts in case composition. We do
\emph{not} interpret them as causal effects. Valid causal
inference via IPTW requires three assumptions, each of which
is strained in this setting:
\begin{enumerate}
    \item \textbf{No unmeasured confounding} (conditional
    exchangeability). All common causes of treatment and
    outcome must be observed. This is implausible over a
    34-year window: economic cycles, SEC enforcement
    intensity, judicial attitudes, and plaintiff bar
    strategies all evolved concurrently with the PSLRA and
    are not captured in the IDB.
    \item \textbf{Positivity} (overlap). Every covariate
    stratum must contain both pre- and post-PSLRA cases. The
    11.6:1 post-to-pre sample imbalance leaves many covariate
    cells empty in the pre-PSLRA era, producing extreme
    propensity scores that require trimming.
    \item \textbf{Correct model specification.} The
    propensity score model must be correctly specified. We
    assess sensitivity to this assumption through the
    four-strategy triangulation described above: agreement
    across the unadjusted, regression-adjusted,
    weighted-regression, and marginal structural
    specifications is treated as a limited robustness check,
    not as proof of correct specification.
\end{enumerate}
Furthermore, some IPTW covariates (circuit filing patterns
and statutory basis mix) may themselves lie on post-PSLRA
pathways, complicating even the compositional
interpretation. We return to that post-treatment-bias
concern in Chapter~\ref{ch:discussion} as part of the
thesis-level interpretation. For these reasons, we describe
the IPTW estimates as ``composition-adjusted'' throughout,
never as ``causally attributable.''

\section{Shared Frailty Models for Circuit-Level Heterogeneity}
\label{sec:frailty_method}

The Cox models in Section~\ref{sec:hazard_regression} treat circuit as a fixed effect,
estimating a separate hazard ratio for each circuit relative
to a reference. This approach captures \emph{observed}
circuit-level variation but cannot account for unobserved
institutional factors---local judicial norms, bench expertise
in securities law, or district-level procedural culture---that
induce correlation among cases within the same circuit
\citep{ruetenbudde2019}.

A shared frailty model introduces a circuit-level random
effect. Let $u_c$ denote the random intercept for circuit $c$,
with $u_c \sim N(0, \theta)$, where $\theta$ is the frailty
variance. The cause-specific hazard becomes:
\begin{equation}
\lambda_k(t \mid \mathbf{X}_i, u_{c(i)}) = \lambda_{0k}(t) \exp(\boldsymbol{\beta}_k^\top \mathbf{X}_i + u_{c(i)})
\end{equation}
where $c(i)$ is the circuit of case $i$. The random intercept
$u_c$ acts as an additive shift on the log-hazard scale:
circuits with $u_c > 0$ have uniformly higher event rates, and
those with $u_c < 0$ have lower rates, after conditioning on
covariates. Equivalently, $\exp(u_c)$ is a multiplicative
frailty on the hazard scale with a log-normal distribution.

The frailty variance $\theta$ quantifies the degree of
unobserved circuit heterogeneity. A large $\theta$ indicates
that circuits differ substantially in ways not captured by
observed covariates; $\theta \approx 0$ implies that
fixed-effect covariates adequately account for between-circuit
variation. We estimate frailty models using
\texttt{coxme::coxme()} with a \texttt{(1|circuit)} random
effect for both settlement and dismissal outcomes.

With only 11 circuits containing $\geq 50$ securities cases,
frailty variance estimates carry inherent uncertainty.
We therefore present frailty results as a sensitivity analysis
and pair them with cluster-robust standard errors (obtained via
\texttt{coxph(\ldots, cluster = circuit)}) as the primary
robustness check for within-circuit correlation. The informative
comparison is whether the PSLRA hazard ratio changes
meaningfully between fixed-effect and random-effect
specifications.

\section{Model Comparison and Performance Metrics}

\subsection{Concordance Index}

Harrell's concordance index (C-index) measures the proportion
of comparable case pairs for which predicted risk ordering
agrees with observed event ordering \citep{harrell1982}:
\begin{equation}
C = \frac{\sum_{i \neq j} I(T_i < T_j, \delta_i \neq 0) \cdot I(\widehat{\text{risk}}_i > \widehat{\text{risk}}_j)}{\sum_{i \neq j} I(T_i < T_j, \delta_i \neq 0)}
\end{equation}
where $\widehat{\text{risk}}_i = \hat{\boldsymbol{\beta}}^\top \mathbf{X}_i$ is the linear predictor (risk score) from the fitted model, $C = 0.5$ is random ordering, and $C = 1.0$ is perfect discrimination.

\subsection{Time-Dependent AUC}

Time-dependent AUC quantifies discrimination at specific
prediction horizons by constructing ROC curves from incident
case sensitivity and dynamic specificity \citep{heagerty2000}.
We report AUC at $t^* \in \{1, 2, 3, 5\}$ years.

\subsection{Integrated Brier Score}

The Brier score at horizon $t$ measures mean squared prediction
error with inverse probability of censoring weighting (IPCW)
\citep{graf1999}. For competing risks, the cause-specific Brier 
score evaluates the predicted cumulative incidence $\hat{F}_k(t \mid \mathbf{X}_i)$ 
against the observed event status:
\begin{equation}
\text{BS}_k(t) = \frac{1}{n} \sum_{i=1}^n \widehat{W}_i(t) \left( Y_{ik}(t) - \hat{F}_k(t \mid \mathbf{X}_i) \right)^2
\end{equation}
where $Y_{ik}(t) = I(T_i \leq t, \delta_i = k)$ is the binary indicator that event $k$ occurred prior to time $t$, and $\widehat{W}_i(t)$ represents the IPCW weights derived from the Kaplan-Meier estimate of the censoring distribution $\hat{G}(t)$. The integrated Brier score (IBS) averages $\text{BS}_k(t)$
over a grid of evaluation times. Lower IBS indicates better
calibration; as a benchmark, the Kaplan-Meier marginal
estimator (which ignores all covariates) provides the naive
reference against which covariate-informed models are
compared.

\subsection{Validation Strategy}

For predictive validation, we partition the 12,866-case
extended estimation sample into a 70\% training set and 30\%
test set, yielding 9,005 training cases and 3,861 test
cases. This row set excludes the 86 cases with missing
origin and the 16 D.C.\ Circuit cases omitted from
circuit-adjusted models. The 70/30 ratio balances
training-set size (needed for stable coefficient estimation
with 18 covariate terms) against a test set large enough for
precise C-index and AUC estimation. The split is stratified
by PSLRA regime to ensure that both the training and test
sets preserve the 8:92 pre-to-post-PSLRA ratio, preventing
the smaller pre-PSLRA cohort from being underrepresented in
either partition. Predictive models (Cox and Fine-Gray) are
fit on training data; all C-index and AUC values reported in
Chapter~\ref{ch:results} are computed on the held-out test
set. Because sampling is carried out separately within the
pre- and post-PSLRA strata, the realized split differs by
one case from exact 70/30 rounding of 12,866. For prediction
only, statutory basis is dropped from
the linear predictor because complete separation in one
statutory-basis level produces infinite coefficients, but
the train/test split is applied to the same 12,866-case row
set. Fine-Gray C-index values are computed independently
using the Fine-Gray model's linear predictor (the
subdistribution hazard score), which differs from the
cause-specific Cox linear predictor. Because
\texttt{survival::concordance()} treats competing events as
censored rather than maintaining them in the risk set as the
Fine-Gray model does, the reported Fine-Gray C-index is an
approximation.

IPTW and frailty models use the \emph{full extended
estimation sample} ($n = 12{,}866$) without train/test
splitting, because their purpose is compositional adjustment
and heterogeneity quantification, not prediction. Their
diagnostics are covariate balance (standardized mean
differences and effective sample size for IPTW) and frailty
variance (for shared frailty models), not C-index or AUC.

\section{Limitations of the Modeling Framework}

All models condition on observed administrative covariates.
Unobserved factors may still influence both the probability
and timing of outcomes, and finer-grained clustering at the
judge level is infeasible because the IDB lacks consistent
judge identifiers. We therefore interpret all hazard ratios
as associations conditional on observed covariates and rely
on the robustness checks above to bound sensitivity. The
substantive implications of these limits are taken up in
Chapter~\ref{ch:discussion}.
